\section*{Question 1}
A total charge $Q$ initially dispersed at infinity, is to be brought and uniformly distributed over the surface of a conducting sphere of radius $R$. 
\begin{enumerate}
    \item[(a)] Express the surface charge density $\sigma$ on the sphere. --- \textit{0.2 pts}
    \item[(b)] Express the electrostatic potential $V(\vec{r})$ on the surface of the sphere. --- \textit{0.2 pts}
    \item[(c)] How much work is required to assemble this configuration? Note that electrostatic energy stored in a surface charge distribution $\sigma(\vec{r})$ with a potential $V(\vec{r})$ is given by $U=\frac{1}{2} \int_{S} \sigma(\vec{r}) V(\vec{r}) da$ where $da$ is the area element and the integral is calculated over the surface $S$. --- \textit{0.6 pts}
\end{enumerate}

\section*{Solution:}

\begin{solution}
(a) The surface charge density $\sigma$ is defined as the total charge $Q$ divided by the surface area of the sphere, which is $4\pi R^2$:
\begin{center}
\fbox{
    $\displaystyle \sigma \;=\;\frac{Q}{4\pi R^2}$
}
\end{center}

\noindent
(b) The electrostatic potential $V(\vec{r})$ on the surface of the sphere is constant and given by the expression for the potential due to a point charge at a distance $R$ from the charge:
\begin{center}
\fbox{
    $\displaystyle V_{S}=\frac{1}{4\pi\varepsilon_0}\,\frac{Q}{R}$
}
\end{center}
\noindent
(c) The work required to assemble the configuration is the same as the electrostatic energy stored in the surface charge distribution:

$$\displaystyle U
=\frac{1}{2}\sigma V_{S} \int_S da
=\frac{1}{2}\sigma V_{S} (4\pi R^2) = \frac{1}{2}V_S Q = \frac{1}{2} Q \left(\frac{1}{4\pi\varepsilon_0}\,\frac{Q}{R}\right) = \frac{Q^2}{8\pi\varepsilon_0\,R}$$

\begin{center}
\fbox{
$\displaystyle W =  \frac{Q^2}{8\pi\varepsilon_0\,R}$
}
\end{center}

\end{solution}




